\documentclass {article}
\usepackage{graphicx, amsmath, amssymb}

\title{Descrizione progetto}
\author{Vito Mellone}
\date{}

\begin{document}

\maketitle

\section{Inizializzazione e cammini minimax}

Per prima cosa voglio risolvere il problema di trovare un cammino 
minimax ottimo tra due nodi \textbf{\textit{u}} e \textbf{\textit{v}}.
Per farlo ho idea di procedere così: 

\begin{enumerate}

\item   Creo una classe "Grafo pesato" in cui inserire gli archi
        tra i nodi e i relativi pesi
\item   Individuo la componente connessa più grande
\item   Per trovare un cammino tra \textbf{\textit{u}} e 
        \textbf{\textit{v}} utilizzo un algoritmo ricorsivo
        che sfrutta una priority queue: il passo base considera 
        tutti i nodi adiacenti a \textbf{\textit{u}} e li inserisce
        in coda dando la priorità ai cammini collegati tramite
        archi con costo inferiore, successivamente considera
        il nodo con priorità più alta: se é \textbf{\textit{v}}
        interrompe il processo e restituisce il costo del cammino, 
        altrimenti reitera considerando anche i nodi adiacenti
        al quello attuale.
\item   Dopo aver trovato un cammino tra \textbf{\textit{u}} e \textbf{\textit{v}}, grazie al meccanismo della priority
        queue so che quello é un cammino minimax ottimo: sia $P_0 = u, x_1, ..., x_{n}, v$
        il cammino trovato e supponiamo che il suo costo sia $M_{0}$,
        se per assurdo esistesse un altro cammino $P_1 = u, y_1,...,y_{m},v$ di costo
        $M_1$ tale che $M_1 < M_0$. In $P_0$ esiste un arco $x_k \to x_{k+1}$
        di costo $M_0$ e inoltre so che $u \to y_{1}$ ha costo $c_{u \to y_{1}} \leq M_1 < M_{0} = c_{x_k \to x_{k+1}} $,
        ma allora nell'esecuzione dell'algoritmo di ricerca l'arco
        $u \to y_1$ avrebbe avuto la precedenza su $x_k \to x_{k+1}$ e quindi
        il cammino $P_1$ sarebbe stato trovato prima di $P_0$
             
\end{enumerate}

Al momento il costo di memorizzazione dei dati é lineare nel numero
degli AS e dei cammini BGP, noto dall'output che il grafo é caratterizzato da
una forte sparsità, quindi valuto la possibilità di creare una
struttura grafo che renda più efficiente l'utilizzo della memoria



\end{document}